% Complete DRAM/storage research draft. Evaluation curves are analytical.
% Compile this directory independently with pdfLaTeX and BibTeX.
\documentclass[sigplan,10pt,anonymous,nonacm]{acmart}
\setkeys{acmart.cls}{balance=false}
\usepackage{subcaption}
\usepackage{tabularx}
\setcopyright{none}
\settopmatter{printacmref=false,printccs=false,printfolios=true}
\acmConference[EuroSys '27]{European Conference on Computer Systems}
 {April 19--23, 2027}{Rabat, Morocco}
\acmYear{2027}
\acmDOI{}
\acmISBN{}
\title{AegisKV: Workflow-Aware KV Checkpointing\texorpdfstring{\\}{ }Across DRAM and Storage}
\setlength{\emergencystretch}{1.5em}
\renewcommand{\bibfont}{\normalfont\normalsize}
\newcounter{designalgorithm}
\captionsetup[subfigure]{font=normalsize,labelfont=normalfont,textfont=normalfont}
\AtBeginDocument{\typeout{AEGISKV-GEOMETRY: width=\the\textwidth;
 height=\the\textheight; columnsep=\the\columnsep; baseline=\the\baselineskip}}
\begin{document}
\fancypagestyle{aegispreview}{\fancyhf{}\renewcommand{\headrulewidth}{0pt}
 \renewcommand{\footrulewidth}{0pt}\fancyfoot[C]{\normalsize\thepage}}
\pagestyle{aegispreview}
\begin{abstract}
Multi-agent applications accumulate reusable context while executing
interdependent LLM requests. A worker failure can lose this state and delay
several branches at once, yet protecting every context in fast memory is
expensive. We propose AegisKV, a workflow-aware manager for KV checkpoints
in shared DRAM and SSD or shared storage. AegisKV values DRAM retention,
storage demotion, and final eviction separately, combining normal reuse
with the workflow cost of recovery. It evaluates requests affected by a
common failure together and preserves checkpoint validity during asynchronous
migration. The design requires no dedicated GPU backup region and leaves
the model's attention semantics unchanged. We connect trace-derived
workflow observations and existing reconstruction microbenchmarks to a
reproducible analytical study and a complete experimental protocol.
In 36 declared scenarios, the placement model improves on LRU in 17 and
ties in 19; these are model outcomes, not measurements of an implemented
system. The study identifies when storage bandwidth, prediction error,
and commit lag limit the benefit. End-to-end fault injection and normal
service measurements remain necessary to establish the proposed system's
performance.
\end{abstract}

\maketitle
\thispagestyle{aegispreview}
\raggedbottom

\section{Introduction}
\label{sec:introduction}

A user of a multi-agent application waits for a completed task: a tested
patch, a synthesized report, or an answer assembled from parallel searches.
The service sees many LLM requests, but their value depends on how their
outputs connect. A delay on a critical branch postpones the final result;
an equally long delay on another branch may be absorbed before the join.
Application-aware serving already exploits such structure during normal
execution~\cite{lin2024parrot,pan2025kvflow,ni2026topas}. Failures expose a
different resource-allocation problem: which reusable state should remain
quickly recoverable when one worker interrupts several requests together?

Long contexts make this question consequential. An agent's successive
calls append reasoning, code, and tool results. Losing GPU-resident KV
can force the service to reconstruct a large history before it resumes.
KV checkpointing avoids this work~\cite{strati2024dejavu,jayakody2026ghostserve},
and load-aware recovery distributes requests among surviving
workers~\cite{cao2026lumen}. However, storing every history in fast memory
competes with other contexts. Hierarchical caches already span DRAM and
storage~\cite{qin2025mooncake,xie2026strata}; the remaining question is how
to allocate that hierarchy according to the effect of recovery on the
complete workflow.

Neither recency nor recovery latency alone resolves this question. A
recent, long history can belong to a branch with enough slack to tolerate
storage access, while a smaller, less recent history is on the critical
path. Conversely, high slack does not make a history worthless: evicting
its last copy can cause repeated reconstruction at each normal invocation.
A common worker failure further couples requests through both the DAG and
shared storage, network, and target GPUs. Summing isolated request penalties
can choose a placement that accelerates one branch without advancing its
join. Finally, a placement change is asynchronous: releasing DRAM before
its storage replacement commits can invalidate the very checkpoint the
policy intends to protect.

We propose AegisKV, an online KV checkpoint manager over shared DRAM and
SSD or shared storage. It separates fast-tier retention from last-copy
preservation. The Tracker estimates remaining workflow time; the KV Master
compares complete migration plans using normal reuse and failure-domain
recovery episodes. The Router selects compatible surviving execution
targets, and versioned manifests and reader leases preserve recovery
validity during demotion. Ordinary inference HBM remains available for
execution; the design adds no dedicated Peer-HBM backup region.

The contributions are threefold. First, we formulate tier placement in
terms of the workflow consequences of both normal reuse and common-domain
failure, distinguishing DRAM demotion from final deletion. Second, we
specify a bounded online controller and an asynchronous checkpoint protocol
that account for physical reservations, commit lag, and shared ownership.
Third, we provide reproducible analytical scenarios and a matched
experimental protocol that isolates dependency information, capacity,
storage performance, and normal-service cost. The model exhibits gains,
ties, and nonzero regret against a feasible oracle; it is not evidence of
an implemented system's speedup. Our intended system claim is that
workflow-aware protection improves the allocation of finite recovery
resources, which the planned fault-injection experiments must test.

% Background and Motivation: two background topics, two motivations,
% followed by a standalone failure-position example introducing AegisKV.
\section{Background and Motivation}
\label{sec:background-motivation}

\subsection{Background}
\label{sec:motivation-background}

\paragraph{Multi-agent workflows.}
Multi-agent workflows are becoming an important form of LLM application,
including coding assistants and research systems~\cite{wu2023autogen,anthropic2025multiagent}.
An agent may have a predefined role, such as a planner or reviewer, or be
instantiated as a subagent to complete a particular subtask. Each agent
may issue several LLM requests, invoke tools, and incorporate their results
before producing its output. These interactions form a workflow with
sequential dependencies, parallel branches, and joins. Multiple requests
using the same role may share a compatible prompt prefix, while successive
tool interactions extend an individual agent's history. Thus, a request,
an agent, and the complete user workflow describe different units of
execution and completion.

\paragraph{Failures at serving scale.}
LLM services distribute requests across many GPU workers, with each model
replica potentially spanning several devices. As deployments grow,
failures become an operational concern even when individual devices are
reliable: an unavailable device can interrupt an entire replica and shift
its work to the surviving fleet~\cite{cao2026lumen}. The interruption loses
both serving capacity and access to the failed worker's GPU-resident KV
cache. Recovery therefore requires more than rerouting a request; it must
also supply the state needed to resume computation. KV replication and
checkpointing systems already address this state-recovery
problem~\cite{strati2024dejavu,jayakody2026ghostserve,cao2026lumen}.
We focus on fail-stop GPU/worker failures with a compatible surviving
replica, retained token history and control state, and surviving backup
storage. Under this scope, the application-level consequences of the
interruption depend on both its workflow position and its recovery cost.

\begin{figure*}[t]
  \centering
  \captionsetup[subfigure]{font=normalsize,labelfont=normalfont,
    textfont=normalfont,justification=centering,singlelinecheck=false}
  \begin{subfigure}[t]{0.32\textwidth}
    \vspace{0pt}
    \begin{minipage}[c][2.25in][c]{\linewidth}
      \centering
      \includegraphics[width=\linewidth]{figures/workflow.pdf}
    \end{minipage}
    \caption{Four-branch workflow}
    \label{fig:motivation-workflow}
  \end{subfigure}\hfill
  \begin{subfigure}[t]{0.32\textwidth}
    \vspace{0pt}
    \begin{minipage}[c][2.25in][c]{\linewidth}
      \centering
      \includegraphics[width=\linewidth]{figures/workflow-gaps.pdf}
    \end{minipage}
    \caption{Workflow-maximum gaps}
    \label{fig:motivation-branch-gap}
  \end{subfigure}\hfill
  \begin{subfigure}[t]{0.32\textwidth}
    \vspace{0pt}
    \begin{minipage}[c][2.25in][c]{\linewidth}
      \centering
      \includegraphics[width=\linewidth]{figures/recovery-cost.pdf}
    \end{minipage}
    \caption{Context and Prefill cost}
    \label{fig:motivation-recompute}
  \end{subfigure}
  \caption{Motivating workflow-aware KV protection.
    (a) Selected Claude Code trace example with four parallel branches;
    A0/A5 denote Main before/after them. Bars preserve observed timing;
    dashed tails assume that all branch outputs are required at the join.
    (b) ECDF of the maximum branch-duration gap in each of 121 workflows
    with temporally inferred fan-in; it is not a slack distribution.
    (c) Measured Prefill latency for Qwen2.5-7B-Instruct-1M on two
    A100 PCIe 40\,GB GPUs (TP2, BF16, batch one), with MInference/DCA
    sparse Prefill above 32K, including the final one-token EXTEND.
    Reuse restores the same $(C-1)$-token KV and computes one new token.
    DRAM and shared-filesystem reads/H2D complete before computation.
    Shared FS is the measured DPC mount, with direct I/O and no resident
    Linux file pages; internal storage caches are uncontrolled.
    Full prefill starts from empty KV with 8K chunks.
    HBM/DRAM show medians and min--max ranges of three timings;
    Shared FS/full prefill each have one timing per length.
    K~=~1,024; M~=~$10^6$.}

  \label{fig:motivation}
  \Description{Left: a four-branch workflow with A2 completing last at
    151.6 seconds. Middle: an ECDF of workflow-maximum branch-duration
    differences with a median of 138 seconds and a 90th percentile of
    1173 seconds. Right: four measured Prefill curves for Qwen2.5-7B-Instruct-1M
    from 8K to 1,010,000 tokens: blue circles for HBM reuse, orange squares
    for serialized DRAM reuse, purple diamonds for serialized shared-filesystem
    reuse, and red triangles for full prefill. Both axes are logarithmic.}
\end{figure*}

\subsection{Motivation}

\paragraph{Workflow completion requires different protection priorities.}
\label{sec:motivation-workflow}
The user waits for the complete task, so we use workflow completion time
(WCT), from task release to completion of all required results, as the
primary objective. Minimizing an individual request's or agent's completion
time alone does not capture the effect of its dependencies.
Figure~\ref{fig:motivation}(\subref{fig:motivation-workflow}) illustrates
this distinction with a four-branch coding workflow. A2 completes last,
151.6\,s after the fork, whereas A1 completes at 60.4\,s. If the following
Main stage needs all four outputs, accelerating A1 alone does not advance
the join, while delaying A2 postpones downstream work. The same distinction
matters for fault protection: branches differ in how much additional delay
they can absorb. We call this tolerance \emph{slack}.

To characterize this heterogeneity beyond one example, we analyze the
public SemiAnalysis CC/Weka dataset collected from Claude Code coding
workflows~\cite{semianalysis2026cctraces}. Its filtered release contains
393 sessions. We retain overlapping subagent groups whose branches
complete and are followed by a Main request after the last branch finishes;
Main may also execute while the branches run. This yields 286 groups in
121 workflows. For each group, we compute the longest minus the shortest
branch duration; each workflow contributes its maximum such difference.
Figure~\ref{fig:motivation}(\subref{fig:motivation-branch-gap}) shows a
median of 138\,s and a 90th percentile of 1,173\,s. These statistics show
branch-duration heterogeneity among workflow maxima. The trace lacks
explicit result-consumption edges, so fan-in is inferred from timing and
the gaps are not a distribution of slack at failure time. The example and
distribution together motivate assigning protection priorities according
to workflow impact, rather than treating all agent requests equally.

\paragraph{Growing contexts make KV reconstruction increasingly costly.}
\label{sec:motivation-long-context}
An agent's later requests often include earlier reasoning, code, tool
outputs, and other agents' results. Without context reduction, this history
grows across interactions. KV caching avoids repeatedly processing that
history during normal execution. If a failure makes the active KV
unavailable and no usable copy survives, recovery must rebuild the missing
state before continuing; longer histories can therefore turn a brief
interruption into substantial additional computation.

To isolate the cost of populating KV,
Figure~\ref{fig:motivation}(\subref{fig:motivation-recompute}) compares
HBM, DRAM, and shared-filesystem reuse with full prefill from empty KV.
We measure Qwen2.5-7B-Instruct-1M~\cite{qwen2025qwen25oneM} on two A100
PCIe 40\,GB GPUs with TP2, BF16 weights and KV, and batch one, in one
allocation. The SGLang 0.5.10rc0 worker uses MInference's vertical/slash
sparse attention~\cite{jiang2024minference} with the published Qwen1M
head configuration and DCA position handling. Above the retained 32K
context threshold, even the final one-token EXTEND uses sparse attention;
shorter contexts use grouped dense DCA. GQA K/V is shared without repeated
head copies; prefix chunks group at most two query heads to bound selector
memory, while the final query batches all heads. All paths use the same
implementation and configuration.

Total context is $C=M+1$. Full prefill processes $M$ tokens in chunks of
at most 8,192, then executes one separate token. Every reuse path uses
that same sparse-generated historical KV and executes the same final
token. HBM has the history resident before timing. DRAM restores all
historical KV from pinned host memory and synchronizes the copy stream
before computation. Shared FS first reads a per-rank KV file from the
cluster's DPC shared filesystem into those pinned buffers, then performs
the same H2D and computation. File I/O, H2D and computation are serialized;
there is no compression or selective loading. At 1,010,000 tokens,
the historical KV contains 57.92\,GB across both ranks.
Every path ends at synchronized model-body completion with all $C$
tokens' KV resident, before vocabulary projection. Model loading and
backup creation are excluded; no failure is injected.

At 1,010,000 tokens, HBM, DRAM, shared-filesystem reuse, and full prefill
take 1.305\,s, 2.663\,s, 10.156\,s, and 571.21\,s, respectively.
KV reuse avoids historical computation even with sparse attention
enabled; HBM additionally avoids the measured restore waits.
Its latency is 51.0\% below serialized DRAM in this configuration.
This comparison does not quantify the benefit over transfer--compute
pipelines. Shared-filesystem latency includes file open/read/close,
H2D and the final forward. Direct I/O and page-residency checks exclude
Linux file-page hits, but DPC internal and server caches are uncontrolled;
these results are not cold-device or identified-SSD measurements.

We use one fixed synthetic text stream. HBM and DRAM each have
one warmup and three timings; points and bars show the median and
min--max range. Full prefill and Shared FS each have one timing per
length, with no file-read warmup. These are timing samples, not independent
workloads. Short forwards and seven operator checks precede timing,
including a 1M-context, 8K-query case. The 64K point additionally has a
complete original all-head sparse model reference; 256K and 1.01M have
no separate full-prefill warmup. We verify the complete checkpoint hashes,
historical KV, final hidden states, all-layer new KV, and greedy outputs.
CUDA events confirm serialized transfer and computation. These checks
validate the matched paths under the same sparse configuration, not
equivalence to dense attention or long-context task quality. Complete-history
HBM is an execution-state reference; it does not add a Peer-HBM tier to
the DRAM/storage design. Worker latency alone does not establish
end-to-end WCT benefits.

\begin{figure}[t]
  \centering
  \includegraphics[width=\linewidth]{figures/recovery-exposure.pdf}
  \caption{Workflow position changes the impact of recovery.
    Four independent single-branch failure counterfactuals for
    Figure~\ref{fig:motivation}(\subref{fig:motivation-workflow}),
    at 29.5\,s after the fork. Hatching is delay absorbed by join slack;
    solid fill is WCT loss. Reconstruction uses the active request's
    64-token hash-block input proxy and an assumed 64-token tail, with
    Qwen3-14B/A100 BF16 static estimates (batch and TP one).
    Other execution and queueing are held fixed; this is not a
    simultaneous four-branch fault experiment.}
  \label{fig:motivation-recovery-exposure}
  \Description{For A1 through A4, reconstruction costs are 35.2, 32.2,
    36.7, and 38.9 seconds, while corresponding WCT losses are 0, 32.2,
    17.6, and 17.0 seconds. A1 absorbs all reconstruction delay;
    A2 has zero slack and exposes all of its delay.}
\end{figure}

\paragraph{Failure position determines the workflow benefit of protection.}
The two motivations must be considered together: expensive reconstruction
does not imply an equally large WCT penalty. Figure~\ref{fig:motivation-recovery-exposure}
places an independent hypothetical failure on each active branch of
Figure~\ref{fig:motivation}(\subref{fig:motivation-workflow}) at the same
29.5\,s point. Under the stated join assumption, the observed completion
times provide offline slack, and static estimates provide the additional
reconstruction delay. A1's 35.2\,s cost fits within its 91.2\,s slack,
so it adds no WCT in this counterfactual. A2 costs only 32.2\,s to rebuild
but has zero slack, exposing the entire delay. A3 and A4 expose 17.6 and
17.0\,s, respectively. Hence, the branch with the cheapest reconstruction
can be the most important to protect. This controlled example illustrates
dependency effects; concurrent failures and shared recovery resources can
also delay other requests and must be evaluated separately.

\label{sec:motivation-goals}
These observations lead to AegisKV: a workflow-aware KV protection system
that allocates limited recovery state according to its effect on WCT.
The design connects request-level delay tolerance to the recovery value
of that request's KV, while accounting for DRAM, storage capacity, and shared I/O
consumed by protection. Existing valid cache copies remain usable for
recovery; additional protection must justify its cost to normal serving.
The next section develops the corresponding protection, placement, and
recovery mechanisms.

\section{AegisKV Design}
\label{sec:design}

AegisKV manages reusable checkpoints in a finite shared DRAM pool and a
larger storage tier. Storage may be remote SSD or a shared filesystem;
its failure domain and performance must be specified separately. GPUs
retain their normal execution KV and workspace. There is no dedicated
Peer-HBM protection region. The design below is an implementable contract;
Section~\ref{sec:evaluation} distinguishes its analytical evaluation from
implementation evidence.

\subsection{Design Insights}
\label{sec:design-insights}

\paragraph{Recovery value depends on dependencies.}
A faster restore matters only to the extent that it advances a required
workflow result. Figure~\ref{fig:motivation-recovery-exposure} illustrates
this distinction for independent branch failures. A common worker failure
can interrupt several branches simultaneously. Their penalties cannot be
added: if two critical branches each incur 10\,s, reducing only one to
2\,s leaves their join 10\,s late. Placement must therefore evaluate the
branches affected by the same failure together.

\paragraph{Demotion and deletion have different costs.}
Moving a valid checkpoint from DRAM to storage usually adds I/O; deleting
its last surviving copy can require reconstruction. These operations
need separate values and safety checks. A large-slack agent can use
storage for recovery and still reuse its history in the next interaction.
Slack is not a reason to prohibit ordinary caching.

\paragraph{A desired placement is not recoverable state.}
Copies take time and compete with serving. A checkpoint becomes usable
only after all required objects and metadata commit. Likewise, memory is
reclaimed only after a replacement location is committed and old readers
drain. The controller must optimize the state it can actually achieve,
including the period when both copies consume capacity.

\subsection{System Overview}
\label{sec:system-overview}

The orchestrator supplies workflow IDs, dependencies, accepted inputs, and
execution progress (Figure~\ref{fig:system-overview}). The Tracker estimates
remaining durations and slack. The Router combines locality with load for
normal dispatch and selects a compatible surviving replica at failure
time. The KV Master values candidate changes to DRAM and storage residency.
Adapters expose sealed KV, resource pressure, and recovery progress;
Connectors execute transfers. A surviving directory records immutable
object identities, replica generations, and checkpoint manifests.

\begin{figure}[t]
 \centering
 \includegraphics[width=\columnwidth]{figures/storage-overview.pdf}
 \caption{AegisKV with shared DRAM and a capacity tier. Both DPs retain
 ordinary execution HBM. The Router may choose either compatible surviving
 DP; recovery reads resolve through the shared directory. Demotion commits
 the storage copy before releasing DRAM. Storage is outside the protected
 compute failure domain. Arrows show logical paths, not measured rates.}
 \label{fig:system-overview}
 \Description{A workflow supplies a Tracker and KV Master. A Router and
 surviving directory connect two inference replicas to shared DRAM and
 SSD or shared storage. There is no Peer-HBM backup region.}
\end{figure}

Objects are sealed token blocks containing one layer's complete local KV
shard. Identity includes model and weight revision, tokenizer/prefix hash,
position convention, layer, TP shard, numerical format, and token interval.
A \emph{checkpoint group} is a legal collection of these objects and
model-specific auxiliary state. Placement can move individual objects,
but valuation accounts for their recovery dependencies. Shared physical
objects consume bytes once, while their value includes all affected owners.
Dense and sparse models use the same interface; a sparse adapter must also
capture selector state. Placement never changes attention selections.

\subsection{Workflow Tracker: Remaining Time and Slack}
\label{sec:workflow-protection}

The residual DAG contains unfinished execution segments. Fork/join
boundaries split a parent so child execution is not counted again as
parent waiting. A GBDT predicts each segment's remaining duration from
role/model, elapsed time, known prompt length, and completed calls and
tools. Entry predictions mask fields unavailable before launch.
The target is remaining time directly, not total time minus elapsed time.
Train/test splits separate complete workflows and time periods.

For estimated durations $d_a$, critical-path analysis~\cite{kelley1959cpm}
computes earliest finish $E_a=d_a+\max_{b\in P(a)}E_b$, horizon
$T=\max_a E_a$, latest finish $F_a=\min_{b\in S(a)}(F_b-d_b)$, and
slack $\sigma_a=\max(0,F_a-E_a)$. Empty predecessor maxima are zero;
sinks have $F_a=T$. A request inherits its segment's margin. Multiple
requests do not each receive an independent copy of that margin.

The Tracker refreshes on arrivals, completions, DAG changes, and policy
windows, in $O(|V|+|E|)$ time after prediction. An undeclared future branch,
missing request mapping, or invalid estimate disables dependency-based
priority for that workflow. The fallback uses request recovery costs and
normal reuse, with zero slack. New DAG information invalidates the previous
snapshot. Published trajectory prediction~\cite{zhang2026heddle} motivates
using execution context; it does not establish wall-clock accuracy here.

\subsection{Recovery-aware Router}
\label{sec:recovery-router}

Normal dispatch checks compatible prefixes and foreground load, retaining
Agent affinity when useful. A protection decision does not invalidate an
existing ordinary cache entry. On a confirmed fail-stop event, the Router
fences the old worker, obtains its surviving request log and manifests,
and chooses a healthy, layout-compatible target with feasible execution
space. Shared backing removes the requirement to bind retries to a
particular Peer location. Targets can still become overloaded; both the
profile and evaluation must include readmission and background service.

Before acknowledging an output token, a surviving ordered log stores its
request ID, position, token ID, and source epoch. Prompt and accepted tool
inputs are also recoverable. Retries reconstruct the acknowledged suffix,
suppress duplicate delivery, and reject stale epochs. Exactly-once tool
side effects remain the orchestrator's responsibility. Model reload is
excluded when a compatible replica already survives; it is included when
that assumption does not hold. A GPU-process kill tests a different
failure domain from loss of a host and its local storage.

\subsection{KV Master: Online Retention and Demotion}
\label{sec:kv-master}

\subsubsection{Valuing a placement change}
\label{sec:placement-value}
Let $P$ describe committed locations and outstanding reservations, and
$e$ a possible episode within a fixed horizon $H$. Episodes include
normal next-use events and failures of named source domains. A domain
failure identifies \emph{all} its interrupted requests. The profile
returns additional segment delays $\delta(P,e)$, including checkpoint lag,
restore/reconstruction, target queues, and contended transfers. A max-plus
DAG pass evaluates workflow completion $T_w(P,e)$ using these delays.
For one isolated interrupted segment this reduces to
$\max(0,R_r(P,e)-\sigma_r)$; it is not generally the sum of such terms.

For a candidate $P'$ the controller minimizes the change
\begin{equation}
 \Delta J(P'\mid P)=A(P',P)+
 \sum_e p_e\sum_w\bigl[T_w(P',e)-T_w(P,e)\bigr].
 \label{eq:value}
\end{equation}
$A$ is the foreground delay caused by the \emph{transition}, including
writes and temporary occupancy; episode costs describe service after
that transition and exclude those same costs. All terms are seconds;
$p_e$ are probabilities or explicitly declared expected episode counts,
never raw failure rates. For a first-failure model, domain rates
$\lambda_f$ give $p_f=(\lambda_f/\Lambda)(1-e^{-\Lambda H})$ with
$\Lambda=\sum_f\lambda_f$. Multiple-failure episodes require an extended
model. Normal-use probabilities come from past Agent activations;
no future prompt or completion is visible to the controller.

Compare DRAM demotion against the actual surviving storage copy, including
writeback if absent. Compare storage deletion against the remaining valid
checkpoint or reconstruction. A group's marginal value is recomputed
against the current preview after each choice, including every owner of
a shared object. It is not copied into every block's score. Equal local
slack values can conceal a common-failure join, and the objective remains
non-additive. The greedy procedure below has no optimality or competitive
ratio guarantee.

\subsubsection{Selecting a bounded online plan}
\label{sec:placement-plan}
Newly sealed state is offered to the shared cache regardless of slack,
subject to byte and queue limits. This is an attempt, not an unlimited
storage promise. Existing compatible copies are reused first. A configured
block-aligned interval $Z$ triggers review as KV grows; arrivals, tool
handoffs, significant DAG/load changes, and capacity pressure also trigger
reviews. A review may keep every location unchanged.

A snapshot supplies legal retain, promote, demote, and discard operations.
To relieve DRAM pressure, repeatedly choose the feasible demotion with
smallest predicted increase in $J$ per byte actually releasable, breaking
ties by stable object ID. Revalue the remaining candidates after each
choice. Compare a proposed admission or promotion against this complete
victim plan, including its write cost. Accept an optional change only if
its upper cost bound is below the keep plan's lower bound by a configured
hysteresis margin. Required capacity relief can proceed without a positive
performance benefit, but still obeys safety constraints.

Storage pressure uses a separate victim set and value: only unowned
checkpoints, superseded versions, or explicitly withdrawn protection
claims can lose their last replica. With no safe victim, reject/coalesce
new writes, preserve the last valid frontier, and expose the longer replay
tail. A tool handoff releases an execution reference but does not erase
the cache or automatically withdraw a checkpoint claim. Active requests
can be demoted to storage; physical I/O leases, not permanent DRAM pins,
protect ongoing readers.

\begin{figure}[t]
 \raggedright
 \refstepcounter{designalgorithm}\label{alg:placement}
 \textbf{Algorithm \thedesignalgorithm: Bounded online tier management}
 \par\smallskip
 \begin{tabularx}{\columnwidth}{@{}r@{\quad}X@{}}
 \toprule
 1 & Snapshot DAG, epochs, locations, leases, and byte ledgers \\
 2 & Reuse committed objects; offer newly sealed dependencies \\
 3 & Form legal candidates within the configured work budget \\
 4 & While a preview exceeds a tier's capacity: \\
 5 & \quad Select least-cost legal relief per releasable byte \\
 6 & \quad Update the preview; revalue affected candidates \\
 7 & Compare keep with the complete admission/migration plan \\
 8 & Reject unsupported or unprofitable optional changes \\
 9 & Recheck versions and reserve destination/transfer space \\
 10 & Enqueue copies; retain old authoritative locations \\
 11 & On completion: validate, commit, redirect, drain, release \\
 12 & Publish actual coverage; retry after relevant state changes \\
 \bottomrule
 \end{tabularx}
 \caption*{A preview does not free memory or advance a checkpoint.}
 \Description{An online controller compares complete legal migration plans,
 reserves resources, and releases old copies only after safe commitment.}
\end{figure}

Profiles are keyed by model/layout, bytes, tail length, topology, target
load, and I/O scheduler. Each records support ranges, uncertainty, expiry,
and whether rates are measured or assumed. Missing or stale profiles
freeze discretionary promotions and use bounded LRU demotion with the
same safety rules. The prototype must measure candidate-work budgets,
profile-query overhead, and the resulting regret; unbounded exhaustive
search is only a small-instance evaluation oracle.

\subsubsection{Committing a checkpoint}
\label{sec:checkpoint}
Copies progress through pending, committed, and retiring generations.
Publication requires all required layers, TP shards, token intervals,
selector state, and a matching surviving token-log prefix. The authoritative
frontier is the latest dependency-complete manifest, not the largest
position seen in any individual copy. An incomplete new manifest leaves
the old one usable. Immutable objects can be shared across manifests.

DRAM demotion follows copy--verify--commit--redirect--drain--free. Storage
commit means the backend's documented completion/persistence guarantee;
a buffered write return is insufficient when the claimed failure model
requires durable media. Readers lease a specific location generation.
A failed copy keeps the old location, releases its failed reservation,
and triggers replanning. Duplicate completions and obsolete epochs cannot
publish a new frontier. Last-copy deletion cannot silently preserve the
old claim of recoverability.

\subsection{Adapter and Connector: Safe Concurrent Execution}
\label{sec:kv-runtime}

Each tier accounts for committed bytes, reservations, retiring buffers,
and metadata. Recovery also reserves bounded GPU workspace and staging
buffers before admission. A lower DRAM target can temporarily be below
physical occupancy while safe demotion drains; it cannot be used as an
instantaneous free-space report. Exhausted storage or transfer queues
cause bounded rejection/backpressure, not hidden allocation.

Connectors maintain separate foreground, demand-restore, and writeback
queues. Positive minimum service shares preserve progress for restore
and required demotion; unused service is borrowable. Background backup
has a configured bandwidth cap. A bandwidth share is not a bound on
foreground latency. Byte-based deficit round robin and bounded transfer
quanta prevent one long context from monopolizing a path. Normal reads
and all policies use the same scheduler in controlled comparisons.

The engine consumes chunks at supported kernel boundaries. Layer-wise
prefetch can overlap remaining replay computation, but historical
recomputation already avoided by reuse provides no overlap window.
Missing inputs stall the dependent work, not every retry. TP ranks retain
matched batch and collective order. Dense or native sparse attention
executes unchanged; batched replay requires demonstrated equivalence to
the original state evolution, otherwise the runtime replays token by token.

These rules give three invariants: pending state is never read as committed;
every published dependency resolves to a surviving compatible generation;
and a physical buffer is never reused with outstanding leases. They
establish a protocol argument, not a claim that concurrency bugs, model
parity, or fault-recovery performance have already been validated.


\section{Evaluation}
\label{sec:evaluation}

We ask whether workflow information improves finite-tier allocation,
where that benefit disappears, and what it costs during normal service.
This version contains \emph{analytical results and an experimental design}.
It does not claim completed end-to-end deployment or failure experiments.
Existing trace observations and worker microbenchmarks in
Section~\ref{sec:background-motivation} retain their distinct evidence levels.

\subsection{Methodology and Reproducible Model}
\label{sec:eval-method}

\paragraph{Metrics and controls.}
For workflow $w$ under policy $p$, define the paired failure penalty
$D_{w,p}=\mathrm{WCT}^{\mathrm{fault}}_{w,p}-\mathrm{WCT}^{\mathrm{healthy}}_{w,p}$.
Each policy has its own matched no-failure run. Comparing policies uses
$D_{w,b}-D_{w,p}$, where positive values favor $p$; absolute healthy and
faulty WCT are also reported. The observation cohort includes directly
interrupted workflows and unaffected workflows sharing recovery resources.
Normal-service metrics include TTFT, inter-token delay, SLO goodput,
GPU-seconds, and per-tier request- and token-hit rates. Transfer bytes,
checkpoint age, management CPU time, and failed admissions expose costs
hidden by WCT alone.

\paragraph{Analytical assumptions.}
Four branches have 4, 8, 12, and 16\,GiB checkpoint groups and remaining
times 12, 10, 5, and 1\,s. Their next-use weights are 0.15, 0.35, 0.60,
and 0.85, normalized within a next-use episode. We assume two surviving
targets, a 16\,GiB DRAM budget, and sufficient finite storage (160\,GiB).
DRAM, storage, shared-network, and per-target H2D rates are respectively
20, 4, 12, and 20\,GiB/s. A common 0.25\,s recovery floor and 0.4\,s tail
replay complete the base profile. These are declared generic parameters,
not measurements or calibration of a named GPU or SSD.

An event calendar schedules 0.25\,GiB chunks through a shared DRAM or
storage reader, a shared network, and target-local H2D engines. Equal
chunks enter in round-robin request order. Ready operations take the
earliest available resource, with stable submission order breaking ties.
Tail replay waits for the request's state and shares target compute.
Transport stages overlap; this model does not claim layer-wise
transfer--compute overlap. Recovered branch times propagate through the
join. Placement completes before the modeled fault, so commit lag and
migration are studied separately rather than silently credited as zero.

The scenario grid has three state-size factors (0.5, 1, 2), four slack
factors (0, 0.3, 1, 3), and three storage rates (1, 4, 12\,GiB/s): 36
configurations, not 36 independent physical repetitions. DRAM scales with
state size. The scalar placement score uses failure-episode probability
0.05 and next-use-episode probability 0.02 over a 60\,s horizon; no-event
probability is the remainder. These illustrative probabilities are not a
production failure-rate estimate. The \emph{AegisKV model} greedily admits
whole groups by marginal score reduction per byte. It evaluates the
static value component, not the full asynchronous online algorithm.

\paragraph{Baselines.}
Storage-only retains all state in the capacity tier. LRU fills DRAM in
a declared most-recent order, A3 to A0; reuse-only prioritizes the assumed
next-use probabilities. The feasible oracle enumerates all 16 group
subsets under the same byte budget and minimizes the same score with true
durations. It is an allocation oracle within this restricted group family,
not a globally optimal online system. All-DRAM is an explicitly
capacity-relaxed latency reference. Recompute is separately available in
the model and required in physical tests. Named research systems receive
no fabricated performance curves.

\subsection{Workflow Outcomes and Capacity}
\label{sec:eval-workflow}

\begin{figure*}[t]
 \centering
 \begin{subfigure}{\columnwidth}
  \includegraphics[width=\linewidth]{figures/workflow-outcomes.pdf}
  \caption{Distribution across 36 configurations}
  \label{fig:eval-wct}
 \end{subfigure}\hfill
 \begin{subfigure}{\columnwidth}
  \includegraphics[width=\linewidth]{figures/capacity-frontier.pdf}
  \caption{Finite DRAM, fixed 40\,GiB state}
  \label{fig:eval-capacity}
 \end{subfigure}
 \caption{Analytical placement outcomes. (a) Each point is a deterministic
 scenario from the declared grid; the ECDF is not workload uncertainty.
 (b) Whole checkpoint groups explain discrete capacity transitions.
 All policies use the same queues, byte budget, and surviving targets.
 No confidence intervals or physical repetitions are implied.}
 \label{fig:eval-outcomes}
 \Description{An empirical distribution over enumerated analytical
 configurations and a capacity sweep compare LRU, reuse, workflow value,
 and a capacity-feasible allocation oracle.}
\end{figure*}

Figure~\ref{fig:eval-outcomes} shows the model's workflow penalties.
The AegisKV model improves on LRU in 17 configurations and ties in 19.
Mean penalties are 9.518\,s and 10.892\,s, respectively; these grid means
do not estimate a production improvement. The oracle's mean is 9.386\,s.
The largest greedy regret in the weighted placement objective is
0.204\,s, demonstrating that marginal admission does not solve the
non-additive problem exactly. LRU and reuse-only coincide under the
chosen recency and reuse order; this is a property of these scenarios,
not equivalence of the policies.

The capacity sweep provides two essential controls. With zero DRAM,
policies have the same storage path. Once all 40\,GiB fits, retention
policies converge. Between these endpoints, group sizes, shared queues,
and join sensitivity determine the result. Physical experiments must
sweep the working-set/DRAM ratio, include a no-contention region, and
retain per-workflow negative differences instead of reporting only a
favorable aggregate.

\subsection{Storage, Commit Lag, and Prediction Error}

\begin{figure*}[t]
 \centering
 \begin{subfigure}{\columnwidth}
  \includegraphics[width=\linewidth]{figures/storage-bottleneck.pdf}
  \caption{Storage path sensitivity}
 \end{subfigure}\hfill
 \begin{subfigure}{\columnwidth}
  \includegraphics[width=\linewidth]{figures/checkpoint-lag.pdf}
  \caption{Checkpoint lag and review interval}
 \end{subfigure}
 \caption{Analytical recovery sensitivity. Left: all-DRAM relaxes the
 16\,GiB limit and is only a reference. Right: the ideal phase is enumerated
 over $N=0,\ldots,Z-1$, with 100 generated tokens/s and 2.5\,ms replay/token;
 an additional commit lag adds acknowledged tokens to replay. Other
 parameters are fixed. Neither bandwidth nor replay rate is measured here.}
 \label{fig:eval-storage}
 \Description{Storage bandwidth and checkpoint lag sweeps expose I/O
 bottlenecks and the difference between a review trigger and commitment.}
\end{figure*}

Faster storage narrows the value of DRAM residency until another shared
resource dominates (Figure~\ref{fig:eval-storage}). This is why the main
comparison requires a tuned, pipelined DRAM/storage baseline rather than
only empty-cache recomputation. A shared filesystem must be labeled by
its actual mount, cache state, and network path; a directory name cannot
establish local NVMe behavior. Both client and server caching can change
the apparent bandwidth.

Review interval $Z$ is not a recovery bound. Under the idealized uniform
phase law, the mean uncommitted tail is $(Z-1)/2$ tokens; copying backlog
adds further acknowledged tokens. The right panel explicitly varies that
lag. Actual experiments sample observed commit frontiers and fail at
identical logical progress, with the same Workflow/Agent IDs. They must
also include failures during initial prefill, pending copies, and demotion.

\begin{figure*}[t]
 \centering
 \begin{subfigure}{\columnwidth}
  \includegraphics[width=\linewidth]{figures/prediction-robustness.pdf}
  \caption{Placement under duration perturbations}
 \end{subfigure}\hfill
 \begin{subfigure}{\columnwidth}
  \includegraphics[width=\linewidth]{figures/healthy-io.pdf}
  \caption{Background I/O interference}
 \end{subfigure}
 \caption{Analytical robustness, with no measured predictor or serving
 latency. Left: alternating signs perturb branch duration estimates;
 outcomes use true durations. Right: an M/D/1 queue with 10\,ms service
 quanta and independent Poisson arrivals; curves vary foreground
 utilization while backup adds load. The queue model is separate from
 the restore scheduler and does not predict measured TTFT.}
 \label{fig:eval-robustness}
 \Description{Prediction perturbations affect which groups are placed;
 a separate queueing model illustrates nonlinear interference near saturation.}
\end{figure*}

Prediction error can change which branch is considered critical. The
perturbation study isolates this mechanism; it does not synthesize an
accuracy CDF for the proposed GBDT. Physical validation uses held-out
workflows and causal landmarks, reports absolute seconds-scale errors,
false-positive slack, fallback frequency, and placement regret, and
compares zero-slack, historical, learned, and oracle predictions. Both
systematic bias and out-of-distribution workloads are necessary controls.

\subsection{Attribution, Normal Service, and Safe Migration}

\begin{figure*}[t]
 \centering
 \includegraphics[width=\textwidth]{figures/policy-ablation.pdf}
 \caption{Analytical ablations, averaged equally over 12 size/slack
 configurations at 4\,GiB/s storage bandwidth. No DAG makes all branches
 equally critical; no reuse removes the normal-use term. Fault penalty
 and next-use restore delay are separate outcomes. Bars have no
 uncertainty intervals because the configurations are enumerated,
 not experimental repeats.}
 \label{fig:eval-ablation}
 \Description{Policy ablations compare workflow penalty with normal
 next-use restoration cost, retaining the tradeoff between the two objectives.}
\end{figure*}

Figure~\ref{fig:eval-ablation} removes dependencies and normal-use value
separately. The required physical comparison holds cache capacities,
checkpoint intervals, chunk sizes, transport scheduling, and failure
handling fixed. It changes only the decision rule. A separate comparison
against ordinary hierarchical caching plus a strong load-aware recovery
policy tests whether the combined system improves on an existing stack.
The design currently uses direct marginal revaluation; rental credit is
an optional maintenance alternative whose incremental value must be
isolated under the same score and migration rules.

Normal service is a first-class outcome. Figure~\ref{fig:eval-robustness}
shows, through a separate queueing example, why a 20\% backup bandwidth
cap cannot imply 20\% latency overhead. At higher utilization, queueing
amplifies contention. Full-system runs must therefore report healthy
WCT and SLO goodput as well as injected-failure improvements, including
long-context bursts, repeated Agent calls, and finite storage exhaustion.
A fault-free pass runs each original DAG once, without manufactured reuse
from repeated DAGs or clipping its initial calls.

\begin{figure*}[t]
 \centering
 \begin{subfigure}{\columnwidth}
  \includegraphics[width=\linewidth]{figures/safe-demotion.pdf}
  \caption{A capacity contraction timeline}
 \end{subfigure}\hfill
 \begin{subfigure}{\columnwidth}
  \includegraphics[width=\linewidth]{figures/common-fault.pdf}
  \caption{A common-failure counterexample}
 \end{subfigure}
 \caption{Protocol and objective illustrations. Left: an assumed
 8\,GiB copy at 2\,GiB/s starts at 1\,s; commit/reader drain completes
 at 5.5\,s. The 16\,GiB target does not free physical DRAM early.
 Right: two critical branches each incur 10\,s recovery; retaining
 each one reduces its recovery to 2\,s. The join advances only when
 both are protected. These are exact constructed examples, not runtime tests.}
 \label{fig:eval-safety}
 \Description{A demotion timeline separates capacity targets from actual
 release; a two-branch example separates summed penalties from workflow delay.}
\end{figure*}

Figure~\ref{fig:eval-safety} makes two testable invariants visible. Capacity
reports must follow physical release, and common-failure evaluation must
follow the DAG. Neither example validates concurrent code. The fault
campaign injects failures before/after reserve, copy completion, commit,
redirect, and lease release, plus full storage, duplicate callbacks, and
stale epochs. It checks exact checkpoint bytes and model-state parity,
recovery completion, duplicate-output suppression, and bounded queues.
Correctness failures are reported alongside latency and never dropped.

\subsection{Physical Experiment Plan and Statistical Analysis}
\label{sec:eval-plan}

\begin{figure}[t]
 \centering
 \includegraphics[width=\columnwidth]{figures/controller-scaling.pdf}
 \caption{Controller cost illustration: two DAG passes per candidate and
 eight profile episodes, assuming 0.2\,$\mu$s per node visit. These times
 exclude locks, profile queries, and candidate construction; they are
 not measurements of the proposed controller.}
 \label{fig:eval-controller}
 \Description{An explicitly assumed traversal cost grows with DAG size
 and the candidate-plan budget.}
\end{figure}

Figure~\ref{fig:eval-controller} reserves the control-cost measurement
axes and shows why bounded candidate work requires evaluation. Even a
simple assumed traversal cost can become material as the DAG grows.
Physical profiling must replace that constant and include locks, profile
queries, and the quality lost by limiting the candidate set.

The planned prototype integrates the controller with pinned SGLang and
Mooncake revisions. A minimum deployment has two failure-separated model
replicas and an accessible backing store; 4/8-replica runs evaluate load
redistribution. Candidate models are Qwen2.5-7B-Instruct-1M, Qwen3-8B, and
MiniCPM4-8B, with exact weight/kernel/layout revisions recorded per cell.
Dense and sparse paths are evaluated separately. Context sweeps cover
32K--1M only where the model, kernel, and physical memory support them;
unsupported cells are recorded, not extrapolated as measurements.

Workloads include an instrumented coding workflow, an explicit fan-out QA
workflow~\cite{zhu2024fanoutqa}, and controlled fork/join and sequential
DAGs. CC/Weka supplies locality/timing distributions, not recoverable
prompts or ground-truth dependency edges. Tests vary DRAM/state ratio,
storage capacity and bandwidth, arrivals at 30/60/85\% of the common
baseline's measured saturation point, and failures at 20/50/80/95\% of a
fixed Agent's logical token progress. One source-domain failure interrupts
all colocated requests; request-only interruption is reported separately.

Each primary cell starts with ten independently launched paired campaigns.
A campaign fixes workflow IDs, prompt and tool data, arrival schedule, and
logical failure event across policies, randomizes run order, and includes
a healthy mate for each policy. Report the number of campaigns, workflows,
and requests separately. Use a campaign-level paired bootstrap for mean
differences; hierarchical workload resampling is a secondary analysis.
Tail quantiles require enough completed workflows and remain descriptive
when support is inadequate. Timeouts remain explicit outcomes: report
completion rates and restricted completion time at the same horizon,
without silently deleting censored workflows. Report all planned cells,
confidence intervals, failures, and zero/negative effects. More seeds do
not manufacture additional independent workloads.

The accompanying protocol fixes per-experiment controls, log schemas,
required baselines, and falsification conditions. Identical CSV schemas
and plotting code accept either analytical or measured data; measured
mode additionally checks raw-source hashes and statistical provenance.
Changing data does not require changing figure geometry, but the text and
captions must be revised to match the resulting evidence.

\section{Related Work}
\label{sec:related}

\paragraph{Fault-tolerant LLM serving.}
D\'ej\`aVu streams KV for disaggregation, memory management, and replicated
state recovery~\cite{strati2024dejavu}. GhostServe protects streaming KV
with host-memory erasure-coded parity~\cite{jayakody2026ghostserve}.
RaidServe addresses resilience and imbalance in tensor-parallel serving
through cyclic placement, hybrid attention, and load-aware
routing~\cite{xu2026raidserve}. LUMEN coordinates checkpoint placement,
interrupted-request distribution, and capacity restoration during model
reload~\cite{cao2026lumen}. These works establish that KV state and
recovery load both matter. AegisKV targets a complementary decision:
which contexts deserve finite fast-tier capacity when their restoration
has different workflow consequences. A convincing comparison must pair
its placement with the same load-aware recovery substrate; generic LRU
alone is insufficient as the final systems baseline.

\paragraph{Hierarchical context caching.}
Mooncake pools KV across memory and storage and schedules around cache
locality~\cite{qin2025mooncake}; its SSD offload supports memory-pressure
spill and storage retrieval~\cite{mooncake2026ssdoffload}. CachedAttention
uses hierarchical caching, layer-wise preload, and asynchronous
saving~\cite{gao2024cachedattention}. Strata addresses fragmented I/O,
delay hits, and cache-aware scheduling in long-context
serving~\cite{xie2026strata}. HCache restores states using intermediate
activations and balances I/O with computation~\cite{gao2025hcache}.
AegisKV does not claim hierarchical KV storage or prefetching as new.
Its proposed addition is failure-domain and dependency-aware valuation,
with separate costs for safe demotion and deletion of the last replica.

\paragraph{Workflow-aware serving and prediction.}
Parrot exposes application dataflow to optimize end-to-end
execution~\cite{lin2024parrot}. KVFlow prioritizes prefixes using future
Agent activation distance~\cite{pan2025kvflow}; Pythia and Cortex exploit
workflow semantics for efficient Agent serving~\cite{yu2026pythia,pagonas2025cortex}.
TOPAS jointly chooses prefix retention and execution using remaining
service paths, reuse, and movement costs~\cite{ni2026topas}. This is close
prior work: adding a critical-path score to a cache is not by itself our
novelty claim. The intended distinction is recovery under a common worker
failure, potentially stale checkpoints, and tier transitions that must
preserve a valid recovery state. Heddle's trajectory
prediction~\cite{zhang2026heddle} also does not validate the accuracy of our
wall-clock estimator. These distinctions must survive the ablations in
Section~\ref{sec:evaluation}.

\paragraph{Placement and replacement mechanisms.}
Critical-path analysis~\cite{kelley1959cpm} and cost/size-based online
caching~\cite{young2002landlord} supply useful ingredients. AegisKV's
state-dependent joint-failure value is non-additive and changes with the
DAG, queues, and replicas. The classical Landlord competitive guarantee
therefore does not transfer to this controller. The small-instance oracle
measures approximation loss without claiming a new bound.

\section{Conclusion}
\label{sec:conclusion}

AegisKV treats KV protection as an online decision about workflow progress
across finite DRAM and storage. It preserves normal reuse, values common
failures through dependencies, and separates desired placement from
committed recoverable state. The analytical study identifies a useful
operating region and clear limits: abundant DRAM, faster storage, inaccurate
predictions, and shared queues can reduce the gain. The provided experiment
and plotting artifacts make those hypotheses directly testable. Establishing
system-level benefits still requires the specified implementation,
correctness campaign, and paired physical measurements.

\paragraph{AI assistance.}
Generative AI assisted the drafting, experimental planning, and artifact
code in this working version. Analytical results are explicitly labeled;
author verification of the final implementation and measurements remains
part of completing the study.

\bibliographystyle{ACM-Reference-Format}
\bibliography{references}
\end{document}
